\section{Data}
\label{sec:data}

\subsection{Source}

The dataset is the Kaggle \emph{Anime Recommendations Database}, containing
approximately $73{,}000$ users, $12{,}000$ titles and $7.8$ million ratings.
Each row is a triple \texttt{(user\_id, anime\_id, rating)} where the rating
is an integer in $[1,10]$, or the sentinel $-1$ meaning the user watched the
title but did not rate it.

The dataset is not redistributed with the repository. A generator produces a
synthetic fixture with the same schema, latent taste clusters and a Pareto
popularity distribution, so that the full pipeline and its tests run without
credentials. All numbers in this report are computed on the fixture; the
qualitative conclusions are what transfer, not the absolute values.

\subsection{From ratings to implicit feedback}

We map explicit ratings to implicit positives:
\begin{equation}
  (u,i) \in \Rset
  \iff
  r_{ui} = -1 \;\;\text{or}\;\; r_{ui} \geq 6 .
  \label{eq:positive}
\end{equation}

Two decisions in Equation~\ref{eq:positive} deserve justification.

\paragraph{The sentinel is a positive.} A user who watched a title without
rating it has revealed consumption, which is exactly the signal BPR models.
Discarding these rows would throw away the majority of the interaction data.

\paragraph{Low ratings are dropped, not treated as negatives.} A rating of $2$
is evidence of \emph{dislike}, which is strictly more information than the
``unobserved'' state BPR assumes for the sampled item $j$. Placing a
known-disliked item into the negative pool would be defensible; placing it
into the positive pool would not. We drop these rows because the sampler
draws negatives from the complement of $\Iset_u^{+}$, and a disliked item
correctly belongs there by default.

\subsection{Core filtering}

Users and items with very few interactions carry almost no collaborative
signal and inflate the catalogue over which ranking metrics are computed. We
apply a $k$-core filter with $k = 5$ on both sides.

The filter is applied \emph{iteratively}. Removing low-activity items reduces
user interaction counts, which can push a previously qualifying user below the
threshold; a single pass would leave the stated invariant violated. The loop
runs until a fixed point, so the output genuinely satisfies
\begin{equation}
  |\Iset_u^{+}| \geq k \;\;\forall u \in \Uset
  \qquad\text{and}\qquad
  |\Uset_i^{+}| \geq k \;\;\forall i \in \Iset .
\end{equation}

\subsection{Temporal ordering and its limitation}

The raw dataset carries \textbf{no timestamps}. This is a real constraint and
we state it plainly rather than inventing one. Point-in-time feature
computation requires an ordering, so we derive a per-user sequence index

\begin{equation}
  t_{ui} = \text{position of } i \text{ within the rows of } u
  \quad\text{(file order)} ,
\end{equation}

and treat it as a weak proxy for chronology.

This proxy is adequate for \emph{enforcing} the leakage discipline --- the
machinery that forbids a feature from reading beyond a cut-off is exercised
and tested regardless of whether the cut-off is a true timestamp. It is
\emph{not} adequate for claiming that the evaluation reproduces a real
production split. On a dataset with genuine timestamps, the same code path
takes the timestamp column instead, and the synthetic fixture carries real
timestamps precisely so that the point-in-time tests exercise that path.
